\subsection*{Criterio de Seguridad Física del Personal o Público}
La emergencia podrá ser declarada CLASE I o CLASE II de acuerdo a lo siguiente:
\begin{itemize}
    \item CLASE I: Situación de alto riesgo potencial pero sin que existan lesiones físicas o sobreexposiciones al público.
    \item CLASE II: Situación de alto riesgo donde exista incertidumbre o la certeza de que algún trabajador o persona del público puedan haber recibido una sobreexposición o una lesión física.
\end{itemize}

\subsection*{Criterio de Seguridad Física de la Fuente}
La emergencia podrá ser declarada TIPO A o TIPO B, de acuerdo a lo siguiente:
\begin{itemize}
    \item TIPO A: La Fuente se encuentra en el equipo (Contenedor o Mangueras), no ha sufrido daños y es posible introducirla al blindaje sin incurrir en sobreexposiciones del personal. También se incluyen situaciones donde la integridad del Contenedor se encuentre en duda.
    \item TIPO B: Se ha perdido el control de la Fuente, ya sea porque ésta se encuentra suelta fuera del equipo, o bien porque estando en el equipo, éste fue robado. Se considerará también TIPO B aquellas situaciones donde exista la posibilidad de que la Fuente haya resultado dañada.
\end{itemize}

\subsection*{Para clasificar la emergencia}
Para clasificar la emergencia consulte la figura 2 y en aquellas situaciones en donde exista duda de cómo clasificarla se considerará como la más grave, y a medida que se vayan acumulando datos se podrá ir reclasificando de una manera más acertada. El orden de gravedad se muestra a continuación:
